%=============================================================================
% W4i23 COMPILED UPDATES --- ALL 17 AGENTS
% DFD Research Programme | Iteration 23 (Wave 4, Iteration 14) | 21 March 2026
% Compiler: Opus 4.6
%
% PARADIGM STATE (i22, CONFIRMED):
%   W'(0) = 0 CONFIRMED. v3.2 corrected. W ~ (2/3)y^{3/2} for small y.
%   Deep-MOND G_eff CONFIRMED. Growth delta ~ a^2. sigma_8 = 0.83 +0.22/-0.08.
%   CMB THIRD PEAK: 15-30% deficit (i22) -> 10-25% to ~60% (i23, range widened).
%   beta = 0 at ~2-3 sigma. SO decisive 2027-2028.
%   mochi_class: minimal PoC $5K/1 week. THIS IS THE HIGHEST PRIORITY.
%   SQMS Phase 0: submission-ready. $50K. >30 sigma. SUBMIT NOW.
%   Lab sector 100% decoupled from cosmological questions.
%   Patent portfolio: 4 ALIVE, 6 NICHE, 5 DEAD. STABLE.
%   SCREENING FUNCTION Sigma(y) = 1/sqrt(1+y) NOT DERIVABLE from
%     mu(x) = x/(1+x) via k-essence action. AUTHOR CLARIFICATION NEEDED.
%   Bullet Cluster: JWST 10:1 minor merger reduces deficit to 1.5-2.5x.
%   Prediction count: 9A, 10B, 7C = 26 total. P27 (BAO-a0) added.
%=============================================================================

\documentclass[11pt]{article}
\usepackage[margin=1in]{geometry}
\usepackage{amsmath,amssymb}
\usepackage{hyperref}
\usepackage{enumitem}
\usepackage{booktabs}
\usepackage{xcolor}
\usepackage{tcolorbox}
\usepackage{longtable}

\newcommand{\ee}[1]{\times 10^{#1}}
\newcommand{\DFD}{\textsc{dfd}}
\newcommand{\LCDM}{$\Lambda$CDM}
\newcommand{\Geff}{G_{\rm eff}}
\newcommand{\aM}{a_0}

\title{W4i23 Compiled Updates\\
\large All 17 Agents --- Screening Function Gap Identified,\\
Bullet Cluster Relief, CMB Uncertainty Widened}
\author{DFD Research Programme --- Opus 4.6 Compiler}
\date{21 March 2026}

\begin{document}
\maketitle

\begin{abstract}
Iteration 23 delivers three major results and one critical correction.
\textbf{(1)}~The screening function $\Sigma(y) = 1/\sqrt{1+y}$ cannot
be derived from the k-essence action with $\mu(x) = x/(1+x)$; the
derived form is $\Sigma = (1+y)^2/[y(3+y)]$, which has qualitatively
different asymptotics.  This is the most important theoretical finding
since the W'(0) correction and requires urgent author clarification.
\textbf{(2)}~JWST data (Finner et al.\ 2025) establishes the Bullet
Cluster as a 10:1 minor merger, reducing the DFD mass deficit from
2.3--3.5$\times$ to 1.5--2.5$\times$.
\textbf{(3)}~Deeper CMB third peak analysis reveals that analytic
estimates span 10--61\% depending on the approximation scheme,
confirming that ONLY a Boltzmann code can resolve this question.
\textbf{(4)}~A correction: on the third-peak scale at recombination,
$g_N/\aM \approx 1.6$ (transition regime, NOT deep-MOND), reducing
the $\Geff$ enhancement from $7\times$ to $1.6\times$.

The Euclid DR1 prediction is revised downward: the scale-dependent
$f\sigma_8$ ratio is $\sim 1.05$--$1.15$ (not 1.1--1.8), detectable
at $\sim 2$--$3\sigma$ with four combined bins.  A pre-registration
paper is recommended for submission by August 2026.

Patent portfolio stable (4 ALIVE, 6 NICHE, 5 DEAD).  SQMS Phase 0
and JLab C75 data request remain the top two action items.
\end{abstract}

\tableofcontents
\newpage

%=============================================================================
\section{TOP 5 FINDINGS ACROSS ALL 17 AGENTS}
%=============================================================================

\begin{tcolorbox}[colback=blue!5!white, colframe=blue!70!black,
  title=\textbf{W4i23 TOP 5 FINDINGS (RANKED BY IMPACT)}]

\begin{enumerate}[label=\textbf{F\arabic*.}]

\item \textbf{SCREENING FUNCTION $\Sigma(y) = 1/\sqrt{1+y}$ IS NOT
  DERIVABLE FROM THE k-ESSENCE ACTION.  THE DERIVED FORM HAS DIFFERENT
  ASYMPTOTICS.  ALL CLOCK PREDICTIONS DEPEND ON RESOLVING THIS.}

  The Patent Team (Agents 6--12) attempted a first-principles derivation
  of the screening function from the DFD k-essence action with
  $\mu(x) = x/(1+x)$.  Three approaches were tried:

  \begin{itemize}
  \item \textbf{Static gradient:}  $\Sigma = 1/\mu(y) = (1+y)/y$.
    Diverges as $1/y$ for $y \to 0$; goes to $1$ for $y \to \infty$.
  \item \textbf{Temporal perturbation response:}  $\Sigma \propto
    (1+y)/\sqrt{y(2+y)}$.  Diverges as $1/\sqrt{y}$ for $y \to 0$.
  \item \textbf{Full radial perturbation:}  $\Sigma = (1+y)^2/[y(3+y)]$.
    Diverges as $1/(3y)$ for $y \to 0$; goes to $1$ for $y \to \infty$.
  \end{itemize}

  Compare with $\Sigma(y) = 1/\sqrt{1+y}$: goes to $1$ for $y \to 0$
  and to $1/\sqrt{y}$ for $y \to \infty$.  These are qualitatively
  different (opposite asymptotic behaviors).

  \textbf{The clock paper's screening function is either:}
  (a)~derived from a different physical mechanism than the AQUAL
  potential, (b)~uses a different definition of $y$, or
  (c)~is phenomenological (not derived).

  \textbf{IMPACT:}  All numerical predictions in the clock papers
  (amplitudes, ratios, 7 falsification criteria) depend on $\Sigma$.
  If the functional form changes, the predictions change quantitatively
  (though not qualitatively --- screening still occurs in strong fields).

  \textbf{URGENCY: HIGHEST.  Author clarification needed before any
  further clock paper submissions.}

  \textbf{Grade: B$-$ (honest derivation; important negative result).}

\item \textbf{BULLET CLUSTER IS A 10:1 MINOR MERGER (JWST 2025).
  DFD MASS DEFICIT REDUCED FROM 2.3--3.5$\times$ TO 1.5--2.5$\times$.}

  Finner et al.\ (arXiv:2512.03150, 2025) used joint JWST+DECam
  lensing to measure the first robust virial masses, finding mass
  ratio $\sim 10:1$.  This resolves a 20-year ambiguity and:
  \begin{itemize}
  \item Reduces projection effect contamination in the sub-cluster
    lensing mass by 30--50\%.
  \item Introduces external field effect (EFE) screening of the
    sub-cluster ($g_{\rm ext} \sim 3\aM$).
  \item Revised DFD deficit: 1.5--2.5$\times$.  With 2~eV neutrinos:
    1.2--2.0$\times$.  Potentially consistent within systematics.
  \end{itemize}

  \textbf{The Bullet Cluster is no longer a clear falsification of DFD.
  It remains a tension at the factor-of-two level.}

  \textbf{Grade: B$-$ (new data solid; DFD analysis approximate).}

\item \textbf{CMB THIRD PEAK: ANALYTIC ESTIMATES SPAN 10--61\% DEFICIT.
  THE ANSWER CANNOT BE DETERMINED WITHOUT A BOLTZMANN CODE.}

  The Cosmology Team (Agents 1--5) explored four mechanisms:
  \begin{enumerate}[label=(\alph*)]
  \item $\psi$-potential baryon loading: $f_\psi \approx 0.08$ (revised
    DOWN from 0.15--0.30 at i22, due to transition-regime correction).
  \item $\psi$-acoustic oscillations: NEGLIGIBLE ($< 0.3$\%).
  \item $\Geff$ enhancement of baryon self-gravity: $\Geff/G_N \approx
    1.6$ on third-peak scale (CORRECTED from $7\times$).
  \item Nonlinear mode coupling: ORDER UNITY on CMB scales, confirming
    perturbation theory fails.
  \end{enumerate}

  Different approximation schemes give wildly different results
  (10--61\% deficit), which means:
  \begin{itemize}
  \item The analytic approach has reached its limit.
  \item ONLY the mochi\_class Boltzmann code can determine the true
    third peak height.
  \item The \$5K/1 week minimal PoC is the SINGLE HIGHEST VALUE
    expenditure in the entire DFD programme.
  \end{itemize}

  \textbf{CORRECTION TO i22:}  On the third-peak scale ($\ell \approx
  810$, $k \approx 0.07$~Mpc$^{-1}$) at recombination, $g_N/\aM
  \approx 1.6$ (transition regime), NOT $\ll 1$ (deep MOND).  The
  $\Geff$ enhancement is $\sim 1.6\times$, not $\sim 7\times$.
  The $f_\psi$ baryon loading parameter is correspondingly reduced.

  \textbf{Grade: C (honest about failure of analytic methods).}

\item \textbf{EUCLID DR1 PREDICTION REVISED: SCALE-DEPENDENT $f\sigma_8$
  RATIO $\sim 1.05$--$1.15$, DETECTABLE AT $\sim 2$--$3\sigma$.
  PRE-REGISTRATION PAPER RECOMMENDED.}

  The Predictions Team (Agents 13--17) refined the Euclid DR1
  prediction.  Both $k$-bins ($0.03$ and $0.15$~Mpc$^{-1}$) at
  $z \sim 1$--$1.8$ are in the deep-MOND regime, reducing the
  expected ratio from the i22 estimate (1.1--1.8) to 1.05--1.15.

  A stronger signal comes from the ABSOLUTE $f\sigma_8$ value, where
  a $\sim 20$\% difference from \LCDM{} is predicted but depends on
  CMB normalization (uncertain due to third peak problem).

  Key insight: DFD's faster growth ($f \approx 2$) compensates for
  lower amplitude ($\sigma_8$), making $f\sigma_8$ at $z \sim 1$
  nearly degenerate with \LCDM.  The degeneracy breaks at higher $z$
  or through scale-dependent effects.

  Three independent Euclid DR1 tests:
  (a) scale-dependent $f\sigma_8$ ratio, (b) $\sigma_8(z)$ trajectory,
  (c) $P(k)$ blue tilt.

  \textbf{Recommendation: 4-page pre-registration paper, submit by
  August 2026.}

  \textbf{Grade: B (revised downward; honest about reduced signal).}

\item \textbf{$S_8$ TENSION: DFD QUALITATIVELY AMELIORATES IT BUT
  QUANTITATIVE PREDICTION REQUIRES mochi\_class.  AXIOM STRUCTURE
  STABLE AT 2+1+1 (OR POSSIBLY 1+1+1).}

  DFD's scale-dependent growth naturally produces lower $\sigma_8$ on
  lensing scales than on CMB scales, qualitatively resolving the $S_8$
  tension.  However, the proper $S_8$ comparison requires computing
  the DFD lensing power spectrum with the same pipeline used for \LCDM,
  which is a mochi\_class task.

  The axiom structure is confirmed at: 2 axioms (A1: $n = e^\psi$,
  A2: $c_1 = c\,e^{-\psi}$) + 1 boundary condition ($S^3$) + 1
  derived theorem ($\mu(x) = x/(1+x)$).  Potentially reducible to
  1 axiom (``refractive gravity'') + 1 boundary condition if A2 is
  considered derivable from A1.

  Prediction count updated: 9A + 10B + 7C = 26 total.  One new
  conjecture P27 (BAO amplitude $\propto \sqrt{\aM}$) added.

  \textbf{Grade: B$-$/C+ (mixed).}

\end{enumerate}
\end{tcolorbox}

%=============================================================================
\section{CORRECTIONS TO PREVIOUS ITERATIONS}
%=============================================================================

\begin{tcolorbox}[colback=red!5!white, colframe=red!70!black,
  title=\textbf{CORRECTIONS}]

\begin{enumerate}

\item \textbf{i22 CORRECTION: $g_N/\aM \approx 1.6$ on the third-peak
  scale at recombination (transition regime), NOT deep-MOND.}

  Previous estimate assumed $g_N \ll \aM$ at $z = 1100$ on the
  third-peak scale.  Explicit calculation gives $g_N = 4\pi G\rhob(z_{\rm rec})
  \delta_b/k \approx 2.0\ee{-10}$~m/s$^2$ and $g_N/\aM \approx 1.6$.

  \textbf{Impact:}  $\Geff/G_N$ reduced from $\sim 7$ (deep-MOND)
  to $\sim 1.6$ (transition).  $f_\psi$ reduced from 0.15--0.30 to
  $\sim 0.08$.  CMB third peak deficit estimate range WIDENED (worse,
  not better).

\item \textbf{i22 CORRECTION: Euclid DR1 $f\sigma_8$ ratio reduced
  from 1.1--1.8 to 1.05--1.15.}

  Both $k$-bins at Euclid redshifts are in the deep-MOND regime;
  the ratio between them is smaller than previously estimated.

\item \textbf{i22 FLAG REMAINS: Nuclear clock amplitude discrepancy
  ($10^{-15}$ vs $10^{-18}$) unresolved.}

  The i22 flag about the nuclear clock amplitude stands.  Additionally,
  the screening function derivation gap (Finding F1) may affect the
  amplitude calculation.  Both issues require author clarification.

\end{enumerate}
\end{tcolorbox}

%=============================================================================
\section{ALL 17 AGENT MANDATES: STATUS}
%=============================================================================

\begin{longtable}{p{0.4cm}p{4cm}p{2cm}p{5cm}}
\toprule
\# & Agent Mandate & Status & Key Result \\
\midrule
\endhead
1 & CMB third peak: $\psi$-field & Complete & $f_\psi \approx 0.08$
  (corrected down) \\
2 & CMB third peak: $\Geff$ at $z=1100$ & Complete & $\Geff/G_N \approx
  1.6$ (transition, not deep-MOND) \\
3 & CMB third peak: $\psi$-acoustic & Complete & NEGLIGIBLE ($<0.3$\%) \\
4 & CMB third peak: nonlinear mode & Complete & ORDER UNITY (pert.\ theory
  fails) \\
5 & $\sigma_8/S_8$ analysis & Complete & Qualitative amelioration;
  quant.\ needs mochi\_class \\
6 & Screening derivation & Complete & $\Sigma(y) \neq 1/\sqrt{1+y}$
  from k-essence.  GAP IDENTIFIED. \\
7 & P5 (GW non-lensing) update & Complete & ET on schedule. No change. \\
8 & P7 ($\psi$-trap) update & Complete & DEW competition; focus on
  accelerator/fusion \\
9 & P9 (Q-ratio) update & Complete & JLab C75 data request STILL
  pending \\
10 & P11 (SQMS) update & Complete & Phase 0 ready.  SUBMIT. \\
11 & P1--P4, P6, P8, P10, P12--P15 & Complete & No status changes.
  4 ALIVE, 6 NICHE, 5 DEAD. \\
12 & Bullet Cluster deep dive & Complete & JWST 10:1 minor merger.
  Deficit 1.5--2.5$\times$. \\
13 & Euclid DR1 preparation & Complete & Ratio 1.05--1.15.
  Pre-reg paper recommended. \\
14 & $f\sigma_8$ vs $\sigma_8(z)$ tests & Complete & Near-degeneracy
  with \LCDM{} at $z \sim 1$ in $f\sigma_8$. \\
15 & Axiom independence & Complete & 2 axioms + 1 BC + 1 theorem.
  Stable. \\
16 & New prediction search & Complete & P27 (BAO amplitude $\propto
  \sqrt{\aM}$) added. \\
17 & Prediction audit & Complete & 9A + 10B + 7C = 26 total. \\
\bottomrule
\end{longtable}

%=============================================================================
\section{PARADIGM STATE UPDATE (i23)}
%=============================================================================

\begin{tcolorbox}[colback=green!5!white, colframe=green!70!black,
  title=\textbf{PARADIGM STATE (i23)}]

\begin{itemize}
\item W'(0) = 0 CONFIRMED by author. v3.2 corrected. $W \sim (2/3)y^{3/2}$
  for small $y$.
\item Linear perturbation theory DEGENERATE at FRW background ($\mu(0) = 0$).
\item Deep-MOND $\Geff = G_N\sqrt{\aM k/(4\pi G\rhob\delta_b a)}$ CORRECT.
\item Growth $\delta \sim a^2$ confirmed.  $\sigma_8 = 0.83^{+0.22}_{-0.08}$.
\item \textbf{CMB THIRD PEAK: 10--61\% deficit (WIDENED from 15--30\%).
  ANALYTIC METHODS EXHAUSTED.  mochi\_class REQUIRED.}
\item \textbf{THIRD-PEAK SCALE: transition regime ($g_N/\aM \approx 1.6$),
  NOT deep-MOND.  CORRECTION to i22.}
\item $P(k)$ shape viable on BAO/LSS scales (ratio 0.9--1.5).
\item BAO amplitude enhanced $\times 1.3$--$2.0$.  Phase unchanged.
\item ISW Cold Spot: $-54^{+22}_{-30}~\mu$K vs $-75 \pm 35~\mu$K observed.
\item $\beta = 0$ at $\sim 2$--$3\sigma$.  SO decisive 2027--2028.
\item \textbf{SCREENING FUNCTION: $\Sigma(y) = 1/\sqrt{1+y}$ NOT DERIVABLE
  from $\mu(x) = x/(1+x)$.  AUTHOR CLARIFICATION REQUIRED.}
\item Clock pre-registration v4 complete.  Nuclear amplitude discrepancy
  ($10^{-15}$ vs $10^{-18}$) and screening function gap both UNRESOLVED.
\item mochi\_class: developer-ready spec.  Minimal PoC: \$5K, 1 week.
  \textbf{HIGHEST PRIORITY EXPENDITURE.}
\item \textbf{Bullet Cluster: 10:1 minor merger (JWST 2025).  DFD deficit
  1.5--2.5$\times$ (improved from 2.3--3.5$\times$).}
\item Patent portfolio: 4 ALIVE (P5, P7, P9, P11) + 6 NICHE + 5 DEAD.
  STABLE.
\item Prediction count: 9A + 10B + 7C = 26 total.  P27 new.
\item Euclid DR1 $f\sigma_8$ ratio: 1.05--1.15 (revised down from
  1.1--1.8).  Pre-reg paper recommended by August 2026.
\end{itemize}

\end{tcolorbox}

%=============================================================================
\section{PRIORITY ACTION ITEMS}
%=============================================================================

\begin{tcolorbox}[colback=yellow!5!white, colframe=yellow!70!black,
  title=\textbf{ACTION ITEMS (RANKED BY PRIORITY)}]

\begin{enumerate}

\item \textbf{CRITICAL --- AUTHOR:}  Clarify screening function
  $\Sigma(y)$.  Is it derived or postulated?  From which functional?
  The k-essence action gives $(1+y)^2/[y(3+y)]$, not $1/\sqrt{1+y}$.
  ALL clock predictions depend on this.

\item \textbf{CRITICAL --- FUND:}  mochi\_class minimal PoC (\$5K,
  1 week).  This resolves the CMB third peak question definitively.
  No further analytic progress is possible.

\item \textbf{HIGH --- SUBMIT:}  SQMS Phase 0 proposal (P11).  \$50K,
  $>30\sigma$ DFD vs BCS discrimination.  Target: Q2--Q3 2026 call.

\item \textbf{HIGH --- SEND:}  JLab C75 HOM data request (P9).  Zero
  cost.  Independent Q-ratio measurement.

\item \textbf{HIGH --- WRITE:}  Euclid DR1 pre-registration paper
  (4 pages).  Submit by August 2026.  Three DFD tests with
  falsification criteria.

\item \textbf{MEDIUM --- RESOLVE:}  Nuclear clock amplitude discrepancy
  ($10^{-15}$ vs $10^{-18}$).  Depends partly on $\Sigma$ resolution.

\item \textbf{MEDIUM --- MONITOR:}  ET site selection (2026).
  Prepare T0 summary for press.

\item \textbf{LOW --- INVESTIGATE:}  $S_8$ in DFD requires specifying
  how lensing surveys infer $\Omega_m$ in a MOND universe.
  mochi\_class prerequisite.

\end{enumerate}
\end{tcolorbox}

%=============================================================================
\section{LOOKING AHEAD: i24 PRIORITIES}
%=============================================================================

\begin{enumerate}
\item \textbf{Screening function resolution:}  If author provides
  clarification, re-derive all clock predictions with corrected $\Sigma$.
  If $\Sigma$ is phenomenological, document this honestly.

\item \textbf{mochi\_class status:}  If funded, track progress.
  First result expected: CMB $C_\ell^{TT}$ from a baryons-only DFD
  universe with AQUAL nonlinearity.

\item \textbf{Euclid DR1 paper:}  Begin drafting.  Specify quantitative
  predictions from the analytic calculations (caveated as approximate).

\item \textbf{DESI Y3 data:}  Expected 2026--2027.  BAO amplitude
  and $f\sigma_8$ measurements.  Prepare DFD predictions for comparison.

\item \textbf{Alternative $\mu$ functions:}  Investigate whether a
  different interpolating function (e.g., $\mu(x) = x/\sqrt{1+x^2}$)
  gives a screening function closer to $1/\sqrt{1+y}$.  This could
  resolve the gap identified at i23.
\end{enumerate}

\end{document}
